Logs written by a deployed policy answer only indirectly how users respond to randomly placed items. We propose two log-only signals. LEAF summarises each user's long-term click rate on tags, authors and music of videos outside the test pool, conditional on exposure, and passes these affinities to a graph recommender. DAX models each logged impression as an examination, depending on the user and scenario, followed by an attraction, depending on the item. On KuaiRand-27K, with random insertions for validation and testing, the combined model reaches an nDCG@10 of 0.4854, against 0.4634 for a propensity-weighted LightGCN (PW-LightGCN), our unpublished baseline. The margins over PW-LightGCN and over DIN with DAX are significant but narrow. The PW-LightGCN margin is unchanged under two further training seeds, and the DIN margin is larger under matched capacity and those seeds, all significant at split level. DAX increases per-user nDCG@10 on all six backbones, and at split level on five. Correcting click rates for examination makes the item score useful: removing the examination model significantly lowers the gain, and alone DAX does no better than popularity. Shuffling lifelong statistics across users, or permuting attractiveness across items, removes the gain. The evidence comes from one platform.
